% Source: https://artofproblemsolving.com/community/c3h2495432_regular_heptagon
% Stable identity: aops-topic-2495432-q01
Each interior angle of a regular heptagon is $\dfrac{(7-2)\cdot180^\circ}{7}=\dfrac{900^\circ}{7}$, and all its sides are equal.

\begin{asy}
usepackage("amsmath");
size(5.5cm);
defaultpen(fontsize(10pt));
real al=360/7;
pair P(int k){return dir(270-1.5*al+k*al);}
pair A=P(0),B=P(1),C=P(2),D=P(3),Ee=P(4),F=P(5),G=P(6);
draw(A--B--C--D--Ee--F--G--cycle,linewidth(0.8));
draw(A--C,dashed+gray); draw(B--D,dashed+gray);
draw(A--Ee,dotted);
draw(A--D,blue+linewidth(1));
draw(arc(A,0.2,degrees(B-A)-360,0),red);
draw(arc(D,0.2,180,degrees(C-D)),red);
label("$y$",A+0.34*dir(-12),red+fontsize(9pt));
label("$x$",D+0.34*dir(192),red+fontsize(9pt));
label("$A$",A,unit(A)); label("$B$",B,unit(B)); label("$C$",C,unit(C)); label("$D$",D,unit(D));
label("$E$",Ee,unit(Ee)); label("$F$",F,unit(F)); label("$G$",G,unit(G));
\end{asy}

\textbf{(a)} Triangles $ABC$ and $DCB$ have $AB=DC$, the common side $BC$, and $\angle ABC=\angle DCB=\dfrac{900^\circ}{7}$. By SAS they are congruent, so the diagonals $AC$ and $DB$ (dashed) are equal.

Now triangles $BAD$ and $CDA$ have $AB=DC$, $BD=CA$ and the common side $AD$. By SSS they are congruent, and the angles opposite the equal sides $BD$ and $CA$ are equal:
\[y=\angle DAB=\angle ADC=x.\]
The quadrilateral $ABCD$ is convex (its vertices are vertices of a convex polygon), so its angles add up to $360^\circ$:
\[x+y+\dfrac{900^\circ}{7}+\dfrac{900^\circ}{7}=360^\circ\implies 2x=\dfrac{720^\circ}{7}\implies x=y=\boxed{\dfrac{360^\circ}{7}}.\]
Then $\angle DAB+\angle ABC=\dfrac{360^\circ}{7}+\dfrac{900^\circ}{7}=180^\circ$, so the co-interior angles on the transversal $AB$ show that $AD\parallel BC$. On the other hand $\angle ABC+\angle BCD=\dfrac{1800^\circ}{7}\ne180^\circ$, so $AB$ and $DC$ are not parallel. Hence $ABCD$ is a trapezoid with bases $AD$ and $BC$, equal legs $AB=DC$ and equal base angles at $A$ and $D$: an isosceles trapezoid.

\textbf{(b)} All the student's statements are true, but substituting $b$ gives
\[2y+180^\circ-\dfrac{1800^\circ}{7}+2x=\dfrac{900^\circ}{7}\iff 2x+2y=\dfrac{1440^\circ}{7}\iff x+y=\dfrac{720^\circ}{7}.\]
This is exactly the angle sum of quadrilateral $ABCD$ used in (a). It is an equation in $x+y$ only, so it is satisfied equally well by pairs with $x\ne y$ (for example $x=\dfrac{360^\circ}{7}+5^\circ$, $y=\dfrac{360^\circ}{7}-5^\circ$), and no algebra can extract $x=y$ from it. The computation only uses angle sums, which treat $x$ and $y$ symmetrically; to separate them one needs a genuinely new fact relating the two angles, such as the congruent triangles in (a).
